% 来源：成文/A\_11\_模型评价与改进.md
\section{模型评价与推广}
\subsection{模型优点}
\textbf{优点 1　离散格式严格守恒，且其数学性质可预先自检。} 空间上采用元体平衡（单元中心有限体积），保证逐控制体严格守恒；含水率方程的离散矩阵严格对角占优、非对角元非负，构成 M 矩阵，配合非负性与次随机性质，使离散极值原理先有理论、后设断言——本模型的物理自检不是事后凑数，而是格式性质的直接推论。其边界为固定计算域（问题一至问题三），问题四的移动边界另见 5.4.5。

\textbf{优点 2　温度场时间方向精确推进，时间离散误差归零。} 针对常系数线性系统这一结构性质，采用矩阵指数与分段线性杜哈梅尔积分做解析推进，使时间维不再引入近似。效果可以量化：步长由 1 s 放宽到 60 s，结果仅相差 3.1×10\textsuperscript{-11} ℃。这意味着该场的精度不再受时间步长支配，而只由空间离散决定，从而把调参的自由度集中到真正决定精度的那一维上。

\textbf{优点 3　中心奇点用两条独立推导互证，装配可靠。} 中心控制体为半格，其系数的两条推导路径，即洛必达极限与半格通量平衡，相互独立，所得系数一致，解级比对差异为 2.8×10\textsuperscript{-14} K。这意味着最容易出错的一处几何奇异点由两条互不依赖的路径交叉确认，而不是依赖单一路径的自洽。

\textbf{优点 4　误差带定量给出，验证体系可迁移。} 报告值均附离散误差带（问题三为 ±0.08 h、问题四为 ±0.0249 h），且两类不确定度分列报告、不得相加；所建立的"收敛性检验、解析对拍、独立实现互验、守恒核算、灵敏度分析"五类手段，在变物性与移动边界情形下同样适用，无需重构验证方案。这一点使本文的验证工作不止服务于本题，而具有可复用的方法学价值。

\textbf{定位说明}：本文为单一机理模型的四问递进求解，未采用多模型融合——本题属完全自包含型（数据、参数与物性公式均由题面给出），引入数据驱动模型缺乏训练数据支撑，故改进集中于数值格式与问题特化策略。

\subsection{模型不足与改进方向}
\textbf{不足 1　物理简化带来系统性但已量化的偏差。} 本文忽略蒸发潜热与热质交叉扩散（对应假设 4），并把表面换热与传质系数全程沿用附录 2 的取值（对应假设 5）、传质驱动力取干基浓度差（对应假设 6）。经对照计算，计入潜热后温度系统性变化 0.07–0.10 ℃、问题四的烘干时长仅变化 0.018\%，属已知且有界。

改进方向：若可获得汽化潜热与吸湿等温线，可把能量方程加回源项、把传质驱动力改为活度差，并用干燥实验标定表面换热与传质系数，从而把本节的三项简化一并解除。

\textbf{不足 2　降维与几何演化假设限制了适用面。} 一维轴对称（对应假设 1）忽略了端面效应，两端面面积约占总表面积 7.4\%，使轴中部结果的偏差方向略偏保守；问题四采用仿射收缩（对应假设 8），未考虑表层先干先缩的非均匀分配。

改进方向：对端面效应可用二维轴对称模型做一次对照，本文已由长径比判据说明其量级；对非均匀收缩可改用按局部含水率加权的收缩分配律，并以灵敏度分析检验其影响。需要说明的是，本文未做非均匀收缩的对照计算，此处只给出改进路径而不宣称其结果。

\textbf{不足 3　低含水率端的外推构成主要不确定性。} 扩散系数经验式在低含水率端的外推对烘干时长的贡献，问题三达 43\%–66\%、问题四达 16.8\%–55.1\%，是全文最大的一项模型不确定性；其概率区间所依据的先验分布属主观设定，须在引用时声明，且不得与离散误差带相加。此外，表面含水率的网格分辨率存在限制，界面变系数的取法口径（沿含水率的积分平均）直接决定核心答案。

改进方向：可用低含水率段的独立实验数据标定扩散系数经验式，收窄外推区间；对表面附近做局部加密，进一步压低表面列的分辨率限制；把界面取法的选择作为显式口径随结论一并声明，使读者在引用核心答案时同时看到其口径前提。

\subsection{模型的推广}
本文模型与求解框架可推广到圆柱或球形物料在时变环境下的非稳态传热传质问题，如中药材与食品热风干燥、农产品冷藏温湿耦合、土壤热湿迁移等；其同构之处在于均以守恒律描述内部迁移、以第三类边界描述表面交换、且环境随时间变化。

迁移时需相应调整的内容为四项：物性经验式的形式与适用范围；边界条件的类型；是否计入几何收缩（若计入则沿用本文的物料坐标与守恒形式）；是否计入蒸发潜热与热质交叉扩散。四项均可在本文框架内替换或增补而无需改变离散与求解结构——因为空间离散、时间推进与验证体系均建立在守恒律之上、与具体物性形式解耦。
